
% ---------- Exercise 10 ----------
\item Let

$$
A = \begin{bmatrix}
    \frac{1}{2} & -\frac{1}{2} \\
    -\frac{1}{2} & \frac{1}{2}
\end{bmatrix}.
$$

Compute $A^2$, $A^3$, $A^4$, and $A^5$.\\

\textbf{Solution:}


\begin{answer}
    $$
        A^2 = A \cdot A = \begin{bmatrix}
            \frac{1}{2} & -\frac{1}{2} \\
            -\frac{1}{2} & \frac{1}{2}
        \end{bmatrix} \cdot \begin{bmatrix}
            \frac{1}{2} & -\frac{1}{2} \\
            -\frac{1}{2} & \frac{1}{2}
        \end{bmatrix} = \begin{bmatrix}
            \frac{1}{2} & -\frac{1}{2} \\
            -\frac{1}{2} & \frac{1}{2}
        \end{bmatrix}.
    $$

    Since $A^2=A$, it's clearer that, by the definition of the matrix power:

    $$
        \forall p \in \mathbb{N^+}, \quad A^p = A^{p-1} \cdot A \implies
    $$

    Therefore:
    $$\begin{aligned}
        \implies A^3 & = \underbrace{A^2}_{A} \cdot A = A \cdot A = A^2 = A \\
        \implies A^4 & = \underbrace{A^3}_{A} \cdot A = A \cdot A = A^2 = A \\
        \implies A^5 & = \underbrace{A^4}_{A} \cdot A = A \cdot A = A^2 = A \\
    \end{aligned}$$
\end{answer} \\

\textbf{Alternative solution: (extra)} \\[0.5em]
Since \\
$$
    \boxed{A^2 = A}
$$ \\
we may avoid the efford of calculte the other powers using induction.\\

This means:

    $$
    \boxed{\forall n\in\mathbb{N^+},\quad A^n=A}
    $$

    \textbf{Base case ($n=1$):}

    $$
    A^1=A.
    $$

    \textbf{Inductive hypothesis:}

    Let $n\in\mathbb{N}_{>0}$ and assume that $A^n=A$.\\

    \textbf{Inductive step:}

    Multiplying both sides of the inductive hypothesis
    by $A$ on the right:

    $$
    A^n=A \implies A^nA=AA.
    $$

    Therefore,

    $$
    \begin{aligned}
        A^{n+1}
        &= A^nA \\
        &= AA \\
        &= A^2 \\
        &= A.
    \end{aligned}
    $$

    Therefore, $P(n)\implies P(n+1) \implies \boxed{A^n = A}$

